\section{Phương pháp và kiến trúc mô hình}

\subsection{Biểu diễn dữ liệu thành chuỗi}\label{sec:repr}

Mỗi dòng của BitcoinHeist là một vector đặc trưng, không có cấu trúc chuỗi sẵn có. Đồ án thử hai cách chuyển đổi (Hình~\ref{fig:repr}).

\begin{figure}[H]\centering
\begin{tikzpicture}[font=\small, tok/.style={draw,rounded corners=1pt,minimum width=0.95cm,minimum height=0.55cm,fill=blue!8},
  hstep/.style={draw,rounded corners=1pt,minimum width=1.25cm,minimum height=0.55cm,fill=orange!12},
  pad/.style={draw,dashed,rounded corners=1pt,minimum width=1.25cm,minimum height=0.55cm,fill=gray!8}]
\node[anchor=west] at (-0.6,1.0) {\textbf{(A) Feature-as-sequence}: 1 dòng $\to$ chuỗi 11 bước, bước $i$ = đặc trưng thứ $i$};
\foreach \n [count=\i from 0] in {length,weight,count,looped,neigh.,income,inc\_tz,inc/cnt,inc/nb,w/cnt,loop/cnt}
  \node[tok] (a\i) at (\i*1.08,0) {\scriptsize \n};
\node[below=0.15cm of a5] {\scriptsize $x_i \;\mapsto\; x_i\mathbf{W}_i+\mathbf{b}_i\in\mathbb{R}^{16}$ \quad (thứ tự đặc trưng do ta đặt, không mang nghĩa thời gian)};
\node[anchor=west] at (-0.6,-1.55) {\textbf{(B) Address history}: dòng hiện tại + tối đa $L-1$ lần xuất hiện \emph{trước đó} của cùng địa chỉ};
\node[hstep] (b0) at (0.1,-2.5) {\scriptsize ngày $t_1$};
\node[hstep,right=0.15cm of b0] (b1) {\scriptsize ngày $t_2$};
\node[hstep,right=0.15cm of b1] (b2) {\scriptsize $\cdots$};
\node[hstep,right=0.15cm of b2,fill=red!15,thick] (b3) {\scriptsize hôm nay};
\node[pad,right=0.15cm of b3] (b4) {\scriptsize đệm};
\node[pad,right=0.15cm of b4] (b5) {\scriptsize đệm};
\node[right=0.2cm of b5,align=left] {\scriptsize mỗi bước: 11 đặc trưng\\\scriptsize + cờ ``lần đầu'' + $\log(1+\Delta\text{ngày})$};
\draw[-{Latex}] (b0.south west) ++(0,-0.15) -- ($(b3.south east)+(0,-0.15)$) node[midway,below] {\scriptsize thời gian (chỉ quá khứ, không nhìn tương lai)};
\end{tikzpicture}
\caption{Hai cách biểu diễn một dòng dữ liệu bảng thành chuỗi để đưa vào RNN/CNN.}\label{fig:repr}
\end{figure}

\paragraph{Biểu diễn A --- ``feature-as-sequence''.} 11 đặc trưng số (đã chuẩn hoá) được xem là chuỗi độ dài cố định 11.
Mỗi giá trị vô hướng $x_i$ được ``nhúng'' thành vector $\mathbf{e}_i=x_i\mathbf{W}_i+\mathbf{b}_i\in\mathbb{R}^{16}$, trong đó
$\mathbf{W}_i,\mathbf{b}_i$ là tham số \emph{riêng} của từng đặc trưng (\emph{feature tokenizer}~[8]) --- nhờ đó mô hình biết
bước $i$ là đặc trưng nào. Không có chuỗi dài ngắn khác nhau nên không cần đệm. 27 chiều one-hot thời gian không đưa vào chuỗi
mà được \emph{ghép vào trước lớp phân loại}. \textbf{Hạn chế}: thứ tự đặc trưng là tuỳ ý, không có ``thời gian'' hay ``vị trí''
thật; giả thiết cục bộ của CNN (đặc trưng kề nhau liên quan nhau) và giả thiết tuần tự của RNN chỉ đúng một phần (các đặc trưng
tỉ lệ được đặt cạnh đặc trưng gốc tương ứng). Do đó ta \emph{không kỳ vọng} RNN/CNN vượt mô hình cây trên biểu diễn này.

\paragraph{Biểu diễn B --- ``address history''.} Với mỗi dòng (địa chỉ $a$, ngày $t$), lấy tối đa $L=16$ lần xuất hiện gần nhất
của $a$ tại các ngày $\le t$, sắp theo thời gian, bước cuối là chính dòng hiện tại. Mỗi bước là vector 13 chiều
(11 đặc trưng + cờ ``lần đầu xuất hiện'' + $\log(1+\text{số ngày từ lần trước})$), chiếu tuyến tính sang $\mathbb{R}^{32}$.
\textbf{Chuỗi có độ dài khác nhau} (1--16): đệm 0 ở cuối, kèm mặt nạ (mask); RNN dùng \code{pack\_padded\_sequence} để bỏ qua
phần đệm, CNN nhân mặt nạ trước mỗi lớp tích chập và dùng pooling có mặt nạ. Đây là chuỗi thời gian \emph{thật}, chỉ dùng thông
tin quá khứ (có thể triển khai thực tế: khi một địa chỉ nhận tiền hôm nay, ta đã biết lịch sử của nó). Phải chia dữ liệu theo địa
chỉ để mô hình không ``nhớ'' địa chỉ đã thấy khi train.

\subsection{Mô hình RNN (LSTM / GRU)}

\begin{figure}[H]\centering
\resizebox{\linewidth}{!}{%
\begin{tikzpicture}[font=\small,node distance=0.45cm, box/.style={draw,rounded corners=2pt,align=center,minimum height=0.8cm,fill=#1}, box/.default=white]
\node[box=gray!10] (in) {Chuỗi đầu vào\\ $T\times$ đặc trưng};
\node[box=blue!8,right=of in] (emb) {Tokenizer (A) / \\ Linear (B)\\ $\to T\times d$};
\node[box=green!10,right=of emb] (rnn) {LSTM / GRU\\ $n$ tầng $\times h$ units\\ (Bi-)};
\node[box=green!5,right=of rnn] (pool) {Biểu diễn chuỗi:\\ hidden cuối /\\ mean / attention};
\node[box=orange!10,right=of pool] (cat) {Ghép one-hot\\ thời gian (27)};
\node[box=red!8,right=of cat] (head) {Dense 64, ReLU\\ Dropout\\ Dense 7, softmax};
\foreach \a/\b in {in/emb,emb/rnn,rnn/pool,pool/cat,cat/head} \draw[-{Latex}] (\a) -- (\b);
\end{tikzpicture}}
\caption{Kiến trúc tổng quát của mô hình RNN.}\label{fig:rnn}
\end{figure}

\begin{itemize}
\item \textbf{Lớp hồi tiếp}: LSTM (cổng quên/vào/ra + ô nhớ $c_t$) hoặc GRU (cổng cập nhật/đặt lại, không có ô nhớ riêng, ít
  tham số hơn $\approx$25\%). Cấu hình gốc: 1 tầng, 64 hidden units, một chiều. Các biến thể: 2 tầng, 128 units, hai chiều
  (Bidirectional), không dropout.
\item \textbf{Biểu diễn cuối của chuỗi}: hidden state cuối của tầng trên cùng (mặc định; với Bi-RNN ghép trạng thái cuối của hai
  chiều), trung bình (mean pooling) trên các bước thật, hoặc attention đơn giản $\alpha_t=\mathrm{softmax}_t(\mathbf{v}^\top
  \mathbf{h}_t)$, $\mathbf{z}=\sum_t\alpha_t\mathbf{h}_t$. Ở biểu diễn B, nhờ \code{pack\_padded\_sequence} hidden state cuối chính là
  trạng thái tại ngày hiện tại.
\item \textbf{Dropout}: 0,2 ở lớp Dense cuối và \emph{giữa các tầng} RNN khi có 2 tầng. \emph{Không dùng recurrent dropout}
  (dropout trên kết nối hồi tiếp) vì cài đặt cuDNN của PyTorch không hỗ trợ; với 1 tầng và chuỗi ngắn, dropout ở đầu ra là đủ.
\item \textbf{Đầu ra}: Dense 7 + softmax; hàm mất mát cross-entropy.
\end{itemize}
Chọn \textbf{cả LSTM và GRU} để so sánh kiến trúc, tốc độ và kết quả (mục~\ref{sec:results}).

\subsection{Mô hình CNN (Conv1D)}

\begin{figure}[H]\centering
\resizebox{\linewidth}{!}{%
\begin{tikzpicture}[font=\small,node distance=0.4cm, box/.style={draw,rounded corners=2pt,align=center,minimum height=0.8cm,fill=#1}, box/.default=white]
\node[box=gray!10] (in) {Chuỗi\\ $T\times d$};
\node[box=blue!8,right=of in] (c1) {Conv1D 32, k3\\ (BN) ReLU\\ (Dropout)};
\node[box=blue!4,right=of c1] (p1) {MaxPool 2\\ $T\to\lceil T/2\rceil$};
\node[box=blue!8,right=of p1] (c2) {Conv1D 64, k3\\ (BN) ReLU\\ (Dropout)};
\node[box=blue!4,right=of c2] (p2) {MaxPool 2};
\node[box=green!8,right=of p2] (fl) {Flatten /\\ Global Avg/Max\\ Pooling};
\node[box=red!8,right=of fl] (head) {+ one-hot 27\\ Dense 64, ReLU\\ Dense 7, softmax};
\foreach \a/\b in {in/c1,c1/p1,p1/c2,c2/p2,p2/fl,fl/head} \draw[-{Latex}] (\a) -- (\b);
\end{tikzpicture}}
\caption{Kiến trúc tổng quát của mô hình CNN 1 chiều.}\label{fig:cnn}
\end{figure}

\begin{itemize}
\item \textbf{Loại}: Conv1D (trượt dọc trục chuỗi, mỗi bước có $d$ kênh = chiều embedding). Conv2D không phù hợp vì dữ liệu không có
  cấu trúc 2 chiều.
\item \textbf{Lớp tích chập}: 2 lớp (32 và 64 filters) mặc định; biến thể 3 lớp (32-64-128), rộng (64-128), kernel 3 hoặc 5,
  padding ``same''. Mỗi filter là một bộ phát hiện mẫu \emph{cục bộ} trên $k$ bước liên tiếp, dùng chung trọng số ở mọi vị trí
  (chia sẻ tham số) --- với biểu diễn A là tổ hợp của $k$ đặc trưng kề nhau, với B là mẫu hành vi trong $k$ lần xuất hiện liên tiếp.
\item \textbf{Pooling}: MaxPooling (biến thể AveragePooling) kích thước 2 sau mỗi lớp khi chuỗi còn dài $\ge 4$:
  giảm một nửa độ dài $\Rightarrow$ giảm số tham số của lớp Dense phía sau và tính toán, đồng thời giúp đặc trưng ít nhạy với
  dịch chuyển nhỏ (bất biến cục bộ).
\item \textbf{Batch Normalization} và \textbf{Dropout} (0,3): bật/tắt để so sánh CNN đơn giản với CNN có điều chuẩn.
\item \textbf{Chuyển sang lớp phân loại}: \emph{Flatten} (giữ thông tin vị trí --- hợp với A vì mỗi vị trí là một đặc trưng cố định)
  hoặc \emph{Global Average/Max Pooling} (bất biến độ dài --- bắt buộc với B vì chuỗi dài ngắn khác nhau).
\item \textbf{Phần fully connected}: ghép one-hot thời gian $\to$ Dense 64 ReLU $\to$ Dropout $\to$ Dense 7 softmax.
\end{itemize}

\subsection{Mô hình nền (baseline)}
\begin{itemize}
\item Từ Bài 1 (cùng cách chia ngẫu nhiên, cùng tiền xử lý): Logistic Regression, KNN, Random Forest, HistGradientBoosting,
  LightGBM, XGBoost và Ensemble.
\item \textbf{MLP} (2 tầng ẩn 128 units, ReLU, Dropout 0,2): cùng đặc trưng nhưng \emph{không} coi là chuỗi --- dùng để tách bạch
  lợi ích của ``cấu trúc chuỗi'' khỏi lợi ích của ``mạng nơ-ron''.
\item \textbf{XGBoost trên cách chia theo địa chỉ} với 3 mức thông tin lịch sử: chỉ dòng hiện tại (\code{XGB\_current}); thêm
  số lần đã xuất hiện + khoảng cách ngày (\code{XGB\_count}); thêm trung bình/max của 11 đặc trưng trên 15 lần xuất hiện trước
  (\code{XGB\_hist}). Đây là đối thủ ``dạng bảng'' công bằng của biểu diễn B.
\end{itemize}
